
(4.8.1), on obtient :
\[
  f_*(j_*(y'))+f_*(f^*(x))\in\Fil^{k+1}K^{\bullet}(X)_{\mathbb Q},
\]
soit d'après 3.6 :
\[
  i_*(y)+x\in\Fil^{k+1}K^{\bullet}(X)_{\mathbb Q}.
\]

Pour montrer la relation ii), il faut montrer que :
\[
  y'-g^*(g_*(y'))\lambda_{-1}(F)\in\Fil^kK^{\bullet}(Y')_{\mathbb Q}.
\]
Posons \(z=y'-g^*(g_*(y'))\lambda_{-1}(F)\). On a :
\[
\begin{aligned}
j_*(z)
  &=j_*(y')-j_*\bigl(g^*(g_*(y'))\lambda_{-1}(F)\bigr)\\
  &=j_*(y')-f^*(i_*(y))\qquad\text{d'après 3.4.}
\end{aligned}
\]
Or la relation i) donne :
\[
  f^*(i_*(y))+f^*(x)\in\Fil^{k+1}K^{\bullet}(X')_{\mathbb Q}.
\]
Compte tenu de (4.8.1), on obtient donc :
\[
  j_*(z)\in\Fil^{k+1}K^{\bullet}(X')_{\mathbb Q},
\]
d'où :
\[
  j^*(j_*(z))=z(1-L)\in\Fil^{k+1}K^{\bullet}(Y')_{\mathbb Q}.
\]

Soit \(\bar z\) l'image de \(z\) dans \(\operatorname{Gr}^{k-1}K^{\bullet}(Y')_{\mathbb Q}\), et \(\overline{(1-L)}\) celle de \((1-L)\) dans \(\operatorname{Gr}^{1}K^{\bullet}(Y')_{\mathbb Q}\). On a donc : \(\overline{z(1-L)}=0\). On peut alors appliquer (VI 5.11) (qui reste valable après tensorisation par \(\mathbb Q\), car sa démonstration s'appuie seulement sur la structure de \(K^{\bullet}(Y')\) comme \(K^{\bullet}(Y)\)-algèbre libre), et on peut donc trouver \(z'\in K^{\bullet}(Y)_{\mathbb Q}\) tel que :
\[
  z-g^*(z')\lambda_{-1}(F)\in\Fil^kK^{\bullet}(Y')_{\mathbb Q}.
\]
En appliquant \(g_*\), on trouve, compte tenu de ce que \(g_*(z)=0\) :
\[
  z'\in\Fil^{k+1-d}K^{\bullet}(Y)_{\mathbb Q},
\]
d'où, puisque \(\lambda_{-1}(F)\in\Fil^{d-1}K^{\bullet}(Y')_{\mathbb Q}\) :
\[
  z\in\Fil^kK^{\bullet}(Y')_{\mathbb Q},
\]
ce qui achève la démonstration.

